\section*{Chapitre \ref{ch:g11:vision-and-brain} --- La vision et le cerveau}

\begin{solution}{exo:g11:vision-and-brain:1}
La lumière venue du champ gauche tombe sur la moitié droite de chaque
rétine~; les fibres de la moitié droite de l'\omterm{def:g11:the-eye:parts}{œil} gauche croisent au
chiasma, celles de l'\omterm{def:g11:the-eye:parts}{œil} droit restent à droite~; ensemble elles
forment la bandelette optique droite, font relais dans le thalamus et
atteignent le \omterm{prop:g11:vision-and-brain:pathway}{cortex visuel primaire} droit.
\end{solution}

\begin{solution}{exo:g11:vision-and-brain:2}
Les deux nerfs optiques s'y rejoignent~; les fibres de la moitié nasale
de chaque rétine passent de l'autre côté, celles de la moitié temporale
non. Au-delà, chaque bandelette porte une moitié du champ visuel.
\end{solution}

\begin{solution}{exo:g11:vision-and-brain:3}
L'aire de la couleur (monde vu en gris), l'aire du mouvement (aucune
perception du mouvement), l'aire des visages (visages non reconnus).
\end{solution}

\begin{solution}{exo:g11:vision-and-brain:4}
La plasticité~: la capacité des connexions neuronales à être
renforcées, affaiblies, formées ou élaguées selon l'usage. La période
critique~: la phase précoce de la vie pendant laquelle l'expérience
façonne le cortex de façon décisive et irréversible.
\end{solution}

\begin{solution}{exo:g11:vision-and-brain:5}
Il se fixe sur les récepteurs de la sérotonine des neurones des aires
visuelles, ce qui perturbe le passage des signaux et produit des
hallucinations sans aucune lumière correspondante.
\end{solution}

\begin{solution}{exo:g11:vision-and-brain:6}
Après le chiasma, à gauche~: la bandelette optique gauche, le relais
thalamique gauche ou le \omterm{prop:g11:vision-and-brain:pathway}{cortex visuel primaire} gauche.
\end{solution}

\begin{solution}{exo:g11:vision-and-brain:7}
Les fibres croisées viennent des moitiés nasales des deux rétines, qui
voient les moitiés externes (temporales) du champ~: le patient perd la
moitié gauche du champ de l'\omterm{def:g11:the-eye:parts}{œil} gauche et la moitié droite du champ de
l'\omterm{def:g11:the-eye:parts}{œil} droit --- une vision en tunnel des deux côtés.
\end{solution}

\begin{solution}{exo:g11:vision-and-brain:8}
Normal~: $12 + 20 + 36 + 20 = 88\%$. Privé~: $2 + 3 + 5 + 30 = 40\%$,
dont 10~\% seulement à titre principal.
\end{solution}

\begin{solution}{exo:g11:vision-and-brain:9}
L'occultation force le cortex à utiliser l'entrée de l'\omterm{def:g11:the-eye:parts}{œil} faible, de
sorte que ses connexions sont conservées et renforcées pendant la
période critique au lieu d'être élaguées. Chez l'adulte, la période est
close~: les connexions ne se réorganisent plus à cette échelle, et
l'\omterm{def:g11:the-eye:parts}{œil} faible ne peut pas reprendre son territoire.
\end{solution}

\begin{solution}{exo:g11:vision-and-brain:10}
Ses \omterm{def:g11:the-eye:photoreceptors}{cônes} et leurs \omterm{def:g11:the-eye:photoreceptors}{opsines} sont normaux~; ce qui est perdu, c'est
l'étage cortical qui interprète leurs rapports. Un enregistrement
électrique de la réponse de la rétine à des lumières colorées serait
normal chez elle et anormal chez une personne daltonienne~; et son
déficit est total et acquis, le leur partiel et inné.
\end{solution}

\begin{solution}{exo:g11:vision-and-brain:11}
Des années de pratique ont renforcé et multiplié les connexions au
service de ses doigts, ce qui a agrandi leur territoire cortical. Sans
pratique, l'aire rétrécirait au fil des ans, comme le montrent les
études sur les jongleurs.
\end{solution}

\begin{solution}{exo:g11:vision-and-brain:12}
Elle voit d'abord de la lumière, des couleurs et des taches mobiles
mais ne peut reconnaître ni formes, ni visages, ni profondeur~: la
rétine fonctionne, mais le cortex n'a jamais bâti, pendant la période
critique, les connexions qui interprètent ses signaux. Elle apprendra à
se servir en partie de la vue, lentement, mais jamais avec l'aisance de
quelqu'un qui a vu dès la naissance.
\end{solution}

\begin{solution}{exo:g11:vision-and-brain:13}
À la fovéa, chaque cône a sa propre \omterm{def:g10:cells-common-unit:cell}{cellule} ganglionnaire, si bien que
le centre du champ envoie bien plus de fibres par degré que la
périphérie, où une centaine de \omterm{def:g11:the-eye:photoreceptors}{bâtonnets} en partagent une~; le cortex
répartit sa surface proportionnellement aux fibres, donc au détail, et
non aux degrés de champ.
\end{solution}

\begin{solution}{exo:g11:vision-and-brain:14}
L'activité de cette aire est la perception de la couleur, que le signal
vienne de l'\omterm{def:g11:the-eye:parts}{œil} ou de la mémoire~; la perception est ce que font les
aires, non ce que l'\omterm{def:g11:the-eye:parts}{œil} livre. Une hallucination, c'est cette activité
déclenchée par une drogue au lieu de l'être par l'intention ou par la
lumière.
\end{solution}

\begin{solution}{exo:g11:vision-and-brain:15}
L'\omterm{def:g11:the-eye:parts}{œil} convertit la lumière en signaux~; c'est le cerveau qui voit. Un
patient aux \omterm{def:g11:the-eye:parts}{yeux} intacts peut perdre la couleur, le mouvement ou les
visages par une lésion corticale~; un chaton aux \omterm{def:g11:the-eye:parts}{yeux} intacts devient
fonctionnellement aveugle de l'un d'eux si le cortex n'est pas autorisé
à le connecter~; et une drogue qui ne touche ni l'\omterm{def:g11:the-eye:parts}{œil} ni la structure
du cortex fait voir ce qui n'est pas là. L'\omterm{def:g11:the-eye:parts}{œil} fournit~; le cerveau
voit.
\end{solution}

\begin{solution}{pb:g11:vision-and-brain:1}
\textbf{1.} L'\omterm{def:g11:the-eye:parts}{œil} gauche ou son nerf optique, avant le chiasma~: seules
les fibres d'un seul \omterm{def:g11:the-eye:parts}{œil} sont touchées.

\textbf{2.} La bandelette optique, le thalamus ou le \omterm{prop:g11:vision-and-brain:pathway}{cortex visuel}
droits~: un déficit partagé par les deux \omterm{def:g11:the-eye:parts}{yeux} pour le même demi-champ
ne peut se situer que là où les fibres des deux \omterm{def:g11:the-eye:parts}{yeux} pour cette moitié
ont été rassemblées, c'est-à-dire après le chiasma.

\textbf{3.} Le chiasma lui-même~: les fibres croisées, venues des
moitiés nasales des deux rétines, qui voient les moitiés temporales du
champ, sont sectionnées.

\textbf{4.} L'aire du mouvement, au-delà du cortex primaire. Des champs
et une acuité intacts montrent que le cortex primaire et tout ce qui le
précède fonctionnent.

\textbf{5.} B, C et D ont des rétines intactes~; celle de A peut l'être
ou non. On vérifie en examinant la rétine à l'ophtalmoscope et en
enregistrant sa réponse électrique à des éclairs.

\textbf{6.} Le cortex droit reçoit la moitié gauche du champ.

\textbf{7.} $2^\circ$ centraux~: \qty{1}{cm}~; une bande de $2^\circ$
à $40^\circ$~: \qty{0.1}{cm}. Rapport 10.

\textbf{8.} La fovéa serre ses \omterm{def:g11:the-eye:photoreceptors}{cônes}, chacun avec sa propre \omterm{def:g10:cells-common-unit:cell}{cellule}
ganglionnaire, si bien que les degrés centraux envoient bien plus de
fibres que la périphérie~; le cortex leur donne proportionnellement
plus de surface.

\textbf{9.} Une petite tache aveugle au centre du champ gauche, très
gênante pour la lecture~; plus en avant, une zone aveugle plus large
dans la périphérie du champ gauche, moins remarquée.

\textbf{10.} En bougeant les \omterm{def:g11:the-eye:parts}{yeux}, le patient amène chaque mot sur une
partie de la rétine dont le cortex est intact, en utilisant le pourtour
de la lésion~; c'est lent, mais possible.

\textbf{11.} $20 + 36 + 20 = 76\%$.

\textbf{12.} $2 + 3 + 5 + 30 = 40\%$, et 10~\% seulement avec l'\omterm{def:g11:the-eye:parts}{œil}
gauche comme entrée principale ou à égalité.

\textbf{13.} Son entrée atteint le cortex, mais ses connexions ont été
affaiblies et élaguées faute d'usage~; les connexions de l'\omterm{def:g11:the-eye:parts}{œil} droit se
sont étendues dans l'espace libéré.

\textbf{14.} Que la réorganisation est bornée à une période critique du
début de la vie~; l'équivalent humain est l'amblyopie, qui ne se
développe que dans l'enfance et n'est plus traitable ensuite.

\textbf{15.} Aucun \omterm{def:g11:the-eye:parts}{œil} n'étant favorisé, il n'y a pas d'accaparement~:
l'histogramme reste à peu près symétrique mais avec moins de neurones
réactifs au total, et la vision du chaton est mauvaise des deux \omterm{def:g11:the-eye:parts}{yeux}
--- une privation sans concurrence affaiblit les deux.

\textbf{16.} Le cerveau supprime l'image de l'\omterm{def:g11:the-eye:parts}{œil} dévié pour éviter la
diplopie~; inutilisées, ses connexions sont élaguées pendant la période
critique et l'\omterm{def:g11:the-eye:parts}{œil} devient définitivement faible. Des lunettes corrigent
l'optique, mais l'\omterm{def:g11:the-eye:parts}{œil} n'est pas utilisé, et l'élagage se poursuit.

\textbf{17.} Couvrir le bon \omterm{def:g11:the-eye:parts}{œil} force le cortex à utiliser les signaux
de l'\omterm{def:g11:the-eye:parts}{œil} faible, ce qui conserve et renforce ses connexions.

\textbf{18.} La période critique se termine vers sept ans~; ensuite le
cortex ne redistribue plus le territoire à cette échelle, et les
connexions perdues de l'\omterm{def:g11:the-eye:parts}{œil} faible ne peuvent plus être rebâties.

\textbf{19.} Chez l'adulte, les connexions sont déjà établies et
stables~; l'accaparement par l'autre \omterm{def:g11:the-eye:parts}{œil} est un phénomène du cortex en
développement, pendant la période critique, non du cortex mûr.

\textbf{20.} A~: l'\omterm{def:g11:the-eye:parts}{œil} ou le nerf optique~; C~: le chiasma~; B~: la
bandelette droite, le thalamus ou le cortex primaire~; D~: l'aire du
mouvement au-delà. C'est la plasticité du cortex, bornée à une période
critique, qui fait que l'occultation de l'enfant marche à quatre ans et
échoue à quinze.
\end{solution}
